\section{Causal-JEPA 的证据：增益来自哪里，不能推出什么}

结果部分最容易被压缩成``加了 masking，数字更高''。高质量阅读必须拆开三层：object-centric representation 是否有用，predictor 与 action interface 是否合适，history masking 是否在同一架构上带来额外收益。只有 OC-JEPA 与 C-JEPA 的同架构对照，才最接近回答 masking 本身的贡献。

\subsection{CLEVRER：counterfactual 问题最能暴露交互捷径}

CLEVRER 表格同时报告平均问题准确率与 counterfactual 指标。阅读时应先固定 encoder 与 architecture，再沿 $|M|$ 比较遮蔽对象数。VideoSAUR 表示下，随着适度增加遮蔽，counterfactual per-question accuracy 的增益明显大于总体准确率增益，符合``移除对象后会怎样''与训练时 latent intervention 的结构对应。

\lecturefigure{slide-028.jpg}{CLEVRER：不同 object mask 数量下的推理与 counterfactual VQA}{CS25 V6 official deck, p.~28}

\teachervoice{00:22:38--00:24:39，Hazel 特别要求比较 no-mask 的 OC-JEPA 与同架构 C-JEPA：object-centric representation 不是充分解释，masking 才是 counterfactual 指标额外提升的关键变量。}

\begin{knowledgebox}{读图：为何不能只看最大数字}
表中 VideoSAUR 与 SAVi 是两种 object-centric encoder，$|M|$ 是被遮蔽对象数。适度 masking 提升 counterfactual reasoning，但 SAVi 下过度遮蔽会掉点，说明信息被删除过多时任务本身变得不可辨识。正确结论是存在机制相关的 sweet spot，而不是 ``mask 越多越因果''。
\end{knowledgebox}

\begin{warningbox}{counterfactual VQA 不等于 causal identification}
问题正确率提高说明模型生成的 object trajectories 更支持特定反事实问答；它不证明模型恢复了数据生成图，也不证明在真实干预分布上无偏。benchmark 的场景、问题模板、对象分解与 reasoning head 共同限制了结论范围。
\end{warningbox}

\subsection{Push-T：token efficiency、architecture 与 masking 要分开}

Push-T 表格揭示 object token 的系统优势：patch baseline 使用 $196\times384$ 的 latent 输入，而 object-centric variants 只需约 $6\times128$。但 OC-DINO-WM 直接替换表示后成功率大幅下降，说明 token 更少并不自动更好；静态 object slot 可能没有显式编码速度，原有 causal predictor 与 action concatenation 也未必适配。

\lecturefigure{slide-029.jpg}{Push-T：patch 与 object token 的输入规模及规划成功率}{CS25 V6 official deck, p.~29}

\[
C_{\text{latent}}\propto KdH,
\qquad
\frac{C_{\text{object}}}{C_{\text{patch}}}
\approx\frac{6\times128}{196\times384}\approx1.0\%.
\]
其中，$K$ 是每个时间步 token 数，$d$ 是 token 维度，$H$ 是 rollout horizon，$C_{\text{latent}}$ 是一次 latent rollout 的近似表示计算与存储量。$1\%$ 是输入特征规模的数量级比较，不等价于端到端 wall-clock 必然精确加速一百倍。

下一页把 action conditioning 与 masking 逐项加回。把 action 作为独立 node、改用 bidirectional masked predictor 后，OC-JEPA 相比不适配的 object baseline 显著恢复；在相同 OC-JEPA 架构上增加 object masking，又获得额外提升。因而完整结果来自三件事配合，而非单一``对象化''操作。

\lecturefigure{slide-030.jpg}{Push-T ablation：action node、predictor 与 object masking 的独立贡献}{CS25 V6 official deck, p.~30}

\teachervoice{00:24:39--00:26:43，老师解释 object slot 单帧未必包含 velocity/acceleration，直接沿用 DINO-WM 的 causal predictor 会掉点；动作节点与 bidirectional predictor 先带来约 15 个百分点，masking 相对 OC baseline 再带来约 28 个百分点。}

\begin{importantbox}{Ablation 的正确因果语言}
在给定实验设置内，可以说：\textbf{同一 OC-JEPA 架构中}，history masking 与更高成功率相关，并且去掉 masking 会显著变差。不能说：任何 object-centric world model 加 masking 都会提升 28 个点；encoder、动作表示、horizon、solver 与数据分布都参与结果。
\end{importantbox}

\subsection{PHYRE：相关性 rollout 为什么看起来也像预测}

本节用 PHYRE 的浮空横杆反例，把表格中的成功率差异还原为可观察的机制错误。OC-JEPA 可能学到``两根杆靠近时通常保持相邻''，却没有学习支撑关系；当关键支撑被移走，它仍延续训练相关性。C-JEPA 在训练中反复被问``缺少这个对象自身信息时，哪些对象足以解释未来''，因而更可能把注意力放在真正施加约束的邻近实体上。

\lecturefigure{slide-031.jpg}{PHYRE：学习相关性与学习动力学会产生不同 rollout}{CS25 V6 official deck, p.~31}

下一页把 rollout 与 attention probing 并列。应先看红框中的物理失败，再看表格/热度中 predictor 依赖了谁：OC-JEPA 关注与目标无关的容器或伴随物体，C-JEPA 更关注支撑杆。attention 只能作为依赖线索，不能直接当成 causal edge，但它能与质性失败形成机制一致的证据链。

\lecturefigure{slide-032.jpg}{PHYRE rollout 与 attention probing：错误对象依赖导致物理违例}{CS25 V6 official deck, p.~32}

\teachervoice{00:26:43--00:28:45，课堂把浮空预测解释为 correlation shortcut：模型看到两个对象常同时出现，就预测它们继续一起，而没有识别支撑关系。attention probe 显示失败模型依赖了无关对象，但老师没有把 attention 权重直接宣称为真因果图。}

\begin{warningbox}{attention is not explanation 的本课版本}
attention 权重说明一次前向计算中哪些 token 参与聚合，不能单独证明必要性、充分性或真实因果作用。更强的验证需要 masking/intervention、反事实性能和 rollout 变化共同支持；本课恰好把这三类证据组合起来，但仍受 object encoder 与 benchmark 限制。
\end{warningbox}

\subsection{``causal'' 的操作性定义}

讲者主动把术语收窄为 temporally directed predictive dependencies：历史变量对未来对象状态具有稳定预测贡献，方向由时间给出。它不同于经典结构因果模型中的 graph identification，也不要求每个 latent slot 都是可干预的真实变量。这个限定是理解题目而不被标题误导的关键。

\lecturefigure{slide-033.jpg}{本课中 causal 的限定：时间定向的预测依赖}{CS25 V6 official deck, p.~33}

\[
\mathcal{N}_t^{(i)}
=\arg\min_{\mathcal{N}\subseteq\mathcal{C}_t}
|\mathcal{N}|
\quad\text{s.t.}\quad
p(z_{t+1}^i\mid\mathcal{N})
=p(z_{t+1}^i\mid\mathcal{C}_t).
\]
其中，$z_{t+1}^i$ 是目标对象下一状态，$\mathcal{C}_t$ 是所有可见 context variables，$\mathcal{N}_t^{(i)}$ 是最小 predictively sufficient set。它可被称为 influence neighborhood，但不保证等于真实 causal parents；隐藏 confounder、下游相关变量与部分可观测性都可能改变这个集合。

\subsection{四项假设与两项刻意不假设}

slide 34 将理论主张的支点写得很清楚：没有同一时刻的 instantaneous causality；不同样本共享 transition mechanism；slot 与对象足够对齐；有限 history 足以预测未来。与此同时，方法不要求 first-order Markov，也不要求没有 confounder。前四项使 masked prediction 可解释，后两项让方法能面对更现实的 latent state，但也放弃了真图可辨识性。

\lecturefigure{slide-034.jpg}{Influence neighborhood：成立条件与不要求的理想假设}{CS25 V6 official deck, p.~34}

\begin{knowledgebox}{四项假设逐项翻译}
\begin{enumerate}
  \item \textbf{No instantaneous causality}：作用通过时间滞后表达，不分析同一时间步循环因果；
  \item \textbf{Shared mechanism}：训练视频遵守相对稳定的动力学规则，例如重力不随样本任意改变；
  \item \textbf{Object alignment}：slot 是连贯对象状态，而非频繁交换、拆分或混合；
  \item \textbf{History sufficiency}：给定窗口含有预测所需信息，遗漏的长时记忆不会主导未来。
\end{enumerate}
\end{knowledgebox}

\subsection{现场问题卡：表示不忠实与真图不可恢复}

完整录像在此处出现了一张官方 deck 未包含的问题卡。它问得比结果表更根本：如果 object-centric representation 不 faithful 会怎样？能否恢复 true causal graph？Hazel 的回答是，小误差仍可能保留有用 inductive bias，严重分解错误会让方法失效；当前方法明确不能恢复真图，因为 confounder 与 latent representation 使 identifiability 不成立。

\videofigure{frame-00-31-25.jpg}{现场问题：object representation 失真时会怎样，能否恢复 true causal graph}{Official Stanford recording, 00:31:25}

\teachervoice{00:30:48--00:32:50，老师给出本课最重要的限制：轻微 slot error 可能仍可工作，严重错误则不行；C-JEPA 不恢复 true causal graph，只学习 masking 下有用的 predictive influence。}

\begin{importantbox}{为何后续论文改名也值得记录}
课堂时点使用 ``Object-Level Latent Interventions'' 与 causal inductive bias 语言；2026 年 5 月后的修订将标题改为 ``Object-Level Latent Masking''，并更谨慎地区分可观测性偏置与 causal identification。本讲义按 4 月 9 日课堂快照讲解，同时用这个修订提醒读者：强术语必须随着理论边界更新。
\end{importantbox}

\subsection{回到三项世界模型合同}

第一部分最后回到开场的三列：object representation 回答``现在什么重要''；masking 让 predictor 学习 predictively sufficient context，回答``世界怎样演化''；action as auxiliary node 回答``世界怎样响应你''。这张总结图的价值是把一个论文架构压回可迁移设计原则，而不是记住某个具体网络框图。

\lecturefigure{slide-035.jpg}{Causal-JEPA 对三项世界模型合同的回答}{CS25 V6 official deck, p.~35}

\teachervoice{00:32:50--00:34:51，Hazel 用 object representation、predictive sufficiency 与 action as auxiliary variable 收束前半场，明确表示、转移和动作响应是三个可分别失败的部件。}

\subsection{本章小结}

Causal-JEPA 的 strongest evidence 是同架构 no-mask ablation、counterfactual 指标与 interaction failure 的一致变化。它的 strongest caveat 也同样明确：object alignment 与 history sufficiency 是前提，influence neighborhood 只是 masking 下的 predictively sufficient set，方法不识别 true causal graph。把增益和边界同时保留，才是对这部分课堂内容的忠实重建。

